\begin{proof}
    First it must be shown that $g^\dagger \circ g \approx_E 1_{\dom(g)}$ and $g \circ g^\dagger \approx_E 1_{\cod(g)}$ for each $g \in \Sigma_{\CatHQC}$.
    Let $g \in \Sigma_{\CatHQC}$.
    There are three cases to consider.
    \begin{enumerate}
    \item Assume that $g \in \Sigma_{\mathbf{Prim}}$.
          Then $g^\dagger = g$ and $n = \dom( g ) = \cod( g )$.
          Then it suffices to show that $g \circ g \approx_E 1_n$.
          There are six cases to consider.
          \begin{enumerate}
          \item Assume that $g = \scalebox{0.65}{\circimg{tikz-cache/dagger__term1.pdf}}$.
                Then the following equation holds.
                \begin{equation*}
                    \scalebox{0.65}{\circimg{tikz-cache/proofs__zinv_1.pdf}}
                    \overset{\HqcPowId}{\approx}
                    \scalebox{0.65}{\circimg{tikz-cache/proofs__zinv_2.pdf}}
                    \overset{\HqcPowId}{\approx}
                    \scalebox{0.65}{\circimg{tikz-cache/proofs__zinv_3.pdf}}
                    \overset{\HqcPowAdd}{\approx}
                    \scalebox{0.65}{\circimg{tikz-cache/proofs__zinv_4.pdf}}
                    \overset{\HqcFreePhase}{\approx}
                    \scalebox{0.65}{\circimg{tikz-cache/proofs__zinv_5.pdf}}
                    \overunderset{\HQCCtrlCong}{\HqcPOrd}{\approx}
                    \scalebox{0.65}{\circimg{tikz-cache/proofs__zinv_6.pdf}}
                    \overset{\HqcCtrlId}{\approx}
                    \scalebox{0.65}{\circimg{tikz-cache/proofs__zinv_7.pdf}}
                \end{equation*}
                Then $g \circ g \approx_E 1_n$ by transitivity.
          \item Assume that $g = \scalebox{0.65}{\circimg{tikz-cache/dagger__term8.pdf}}$.
                Then the following equation holds.
                \begin{equation*}
                   \scalebox{0.65}{\circimg{tikz-cache/proofs__pinv_1.pdf}}
                    \overset{\HqcPowId}{\approx}
                    \scalebox{0.65}{\circimg{tikz-cache/proofs__pinv_2.pdf}}
                    \overset{\HqcPowId}{\approx}
                    \scalebox{0.65}{\circimg{tikz-cache/proofs__pinv_3.pdf}}
                    \overset{\HqcPowAdd}{\approx}
                    \scalebox{0.65}{\circimg{tikz-cache/proofs__pinv_4.pdf}}
                    \overset{\HqcPOrd}{\approx}
                    1_0
                \end{equation*}
                Then $g \circ g \approx_E 1_n$ by transitivity.
          \item Assume that $g = \scalebox{0.65}{\circimg{tikz-cache/dagger__term5.pdf}}$.
                Then the following equation holds.
                \begin{align*}
                    \scalebox{0.65}{\circimg{tikz-cache/proofs__hinv_1.pdf}}
                    \;\;\;\overset{\HqcHDef}{\approx}\;\;\;&
                    \scalebox{0.65}{\circimg{tikz-cache/proofs__hinv_2.pdf}}
                    \\
                    \;\;\;\overset{\HqcHDef}{\approx}\;\;\;&
                    \scalebox{0.65}{\circimg{tikz-cache/proofs__hinv_3.pdf}}
                    \\
                    \overset{\text{Lem(\ref{Lem:InvPow})}}{\approx}&
                    \scalebox{0.65}{\circimg{tikz-cache/proofs__hinv_5.pdf}}
                    \\
                    \overset{\text{Lem(\ref{Lem:InvPow})}}{\approx}&
                    \scalebox{0.65}{\circimg{tikz-cache/proofs__hinv_7.pdf}}
                    \\
                    \overset{\text{Lem(\ref{Lem:InvPow})}}{\approx}&
                    \scalebox{0.65}{\circimg{tikz-cache/proofs__hinv_9.pdf}}
                    \\
                    \quad\overset{\text{(a)}}{\approx}\quad&
                    \scalebox{0.65}{\circimg{tikz-cache/proofs__hinv_10.pdf}}
                    \\
                    \overset{\text{Lem(\ref{Lem:InvPow})}}{\approx}&
                    \scalebox{0.65}{\circimg{tikz-cache/proofs__hinv_12.pdf}}
                    \\
                    \overset{\text{Lem(\ref{Lem:InvPow})}}{\approx}&
                    \scalebox{0.65}{\circimg{tikz-cache/proofs__hinv_14.pdf}}
                    \\
                    \overset{\text{Lem(\ref{Lem:InvPow})}}{\approx}&
                    \scalebox{0.65}{\circimg{tikz-cache/proofs__hinv_16.pdf}}
                \end{align*}
                Then $g \circ g \approx_E 1_n$ by transitivity.
          \item Assume that $g = \scalebox{0.65}{\circimg{tikz-cache/dagger__term3.pdf}}$.
                Then the following equation holds.
                \begin{align*}
                    \scalebox{0.62}{\circimg{tikz-cache/proofs__plus_inv_1.pdf}}
                    \overset{\HqcTargDef}{\approx}
                    \scalebox{0.62}{\circimg{tikz-cache/proofs__plus_inv_2.pdf}}
                    \overset{\HqcTargDef}{\approx}
                    \scalebox{0.62}{\circimg{tikz-cache/proofs__plus_inv_3.pdf}}
                    \overset{\text{(c)}}{\approx}
                    \scalebox{0.62}{\circimg{tikz-cache/proofs__plus_inv_4.pdf}}
                    \overset{\text{(a)}}{\approx}
                    \scalebox{0.62}{\circimg{tikz-cache/proofs__plus_inv_5.pdf}}
                    \overset{\text{(c)}}{\approx}
                    \scalebox{0.62}{\circimg{tikz-cache/proofs__plus_inv_6.pdf}}
                \end{align*}
          \item Assume that $g = \scalebox{0.65}{\circimg{tikz-cache/dagger__term2.pdf}}$.
                Then the following equation holds.
                \begin{equation*}
                    \scalebox{0.65}{\circimg{tikz-cache/proofs__ophase_inv_1.pdf}}
                    \overset{\HqcOPhaseDef}{\approx}
                    \scalebox{0.65}{\circimg{tikz-cache/proofs__ophase_inv_2.pdf}}
                    \overset{\HqcOPhaseDef}{\approx}
                    \scalebox{0.65}{\circimg{tikz-cache/proofs__ophase_inv_3.pdf}}
                    \overset{\text{(d)}}{\approx}
                    \scalebox{0.65}{\circimg{tikz-cache/proofs__ophase_inv_4.pdf}}
                    \overset{\text{(a)}}{\approx}
                    \scalebox{0.65}{\circimg{tikz-cache/proofs__ophase_inv_5.pdf}}
                    \overset{\text{(d)}}{\approx}
                    \scalebox{0.65}{\circimg{tikz-cache/proofs__ophase_inv_6.pdf}}
                \end{equation*}
                Then $g \circ g \approx_E 1_n$ by transitivity.
          \item Assume that $g = \scalebox{0.65}{\circimg{tikz-cache/dagger__term4.pdf}}$.
                Then the following equation holds.
                \begin{align*}
                    \scalebox{0.65}{\circimg{tikz-cache/proofs__minus_inv_1.pdf}}
                    \overset{\HqcTargMDef}{\approx}
                    \scalebox{0.65}{\circimg{tikz-cache/proofs__minus_inv_2.pdf}}
                    %\\
                    \overset{\HqcTargMDef}{\approx}
                    \scalebox{0.65}{\circimg{tikz-cache/proofs__minus_inv_3.pdf}}
                    %\\
                    \overset{\text{(a)}}{\approx}
                    \scalebox{0.65}{\circimg{tikz-cache/proofs__minus_inv_4.pdf}}
                    %\\
                    \overset{\text{(d)}}{\approx}
                    \scalebox{0.65}{\circimg{tikz-cache/proofs__minus_inv_5.pdf}}
                    %\\
                    \overset{\text{(a)}}{\approx}
                    \scalebox{0.65}{\circimg{tikz-cache/proofs__minus_inv_6.pdf}}
                \end{align*}
          \end{enumerate}
          In each case $g \circ g \approx_E 1_n$.
          Then  $g^\dagger \circ g \approx_E 1_{\dom(g)}$ and $g \circ g^\dagger \approx_E 1_{\cod(g)}$.
    \item Assume $g \in \Sigma_{\mathbf{Ctrl}}$.
          Then there exists $U \in \Mor(\CatHQC)$ and $V \in \Mor(\CatHQC)$ such that $g = U \odot V$.
          Then the following equations hold.
          \begin{align*}
              g \circ g^\dagger
              &=
              \scalebox{0.65}{\circimg{tikz-cache/proofs__ctrl_dag_b_1.pdf}}
              \overset{\HqcPowId}{\approx}
              \scalebox{0.65}{\circimg{tikz-cache/proofs__ctrl_dag_b_2.pdf}}
              \overset{\text{Lem(\ref{Lem:InvPow})}}{\approx}
              \scalebox{0.65}{\circimg{tikz-cache/proofs__ctrl_dag_3.pdf}}
              \\
              g^\dagger \circ g
              &=
              \scalebox{0.65}{\circimg{tikz-cache/proofs__ctrl_dag_a_1.pdf}}
              \overset{\HqcPowId}{\approx}
              \scalebox{0.65}{\circimg{tikz-cache/proofs__ctrl_dag_a_2.pdf}}
              \overset{\text{Lem(\ref{Lem:InvPow})}}{\approx}
              \scalebox{0.65}{\circimg{tikz-cache/proofs__ctrl_dag_3.pdf}}
          \end{align*}
          Then $g^\dagger \circ g \approx_E 1_{\dom(g)}$ and $g \circ g^\dagger \approx_E 1_{\cod(g)}$ by transitivity.
    \item Assume $g \in \Sigma_{\mathbf{Pow}}$.
          Then there exists $U \in \Mor( \CatHQC )$ and $\alpha \in \mathbb{R}$ such that $g = U \exparr \alpha$.
          It follows that $g^\dagger = U \exparr (-\alpha)$.
          Then $g^\dagger \circ g \approx_E 1_{\dom(g)}$ and $g \circ g^\dagger \approx_E 1_{\cod(g)}$ by \cref{Lem:InvPow}.
    \end{enumerate}
    In each case, $g^\dagger \circ g \approx_E 1_{\dom(g)}$ and $g \circ g^\dagger \approx_E 1_{\cod(g)}$.
    Since $g$ was an arbitrary generator, then $g^\dagger \circ g \approx_E 1_{\dom(g)}$ and $g \circ g^\dagger \approx_E 1_{\cod(g)}$ for each $g \in \Sigma_{\CatHQC}$.
    Then by \cref{Thm:PropInv}, there exists a unique functor $\overline{\dagger}: \CatHQC / E \to \CatHQC / E$ such that $\overline{\dagger} \circ \pi_E = \pi_E^{\text{op}} \circ \dagger$ with every morphism in $\CatHQC / E$ unitary with respect to $( \overline{\dagger} )$.
\end{proof}